\scanpage{151}
\mspage{478}
% position: upper paragraph
D'ailleurs, comme revers de l'aisance dont je parle, j'ai tendance souvent
à sous-estimer la difficulté que peut présenter pour tel élève l'acquisition de
tel bagage, ou le développement de tel outil - ce qui a tendance à le placer en
porte-à-faux vis à vis de mes expectatives. (Voir à ce sujet la note ``Echec d'un
enseignement (1)'', n° 23 iv.) Une telle situation doit assez souvent être un des
ingrédients importants d'une relation fausse au père...

% position: middle note 109
\(({}^{109})\) (9 octobre) Je me suis senti tout content, en terminant la note précé-\linebreak
dente (*), il y a quatre jours. Je me suis trouvé inopinément renouer avec une
intuition qui m'était venue un certain dimanche 17 octobre 1976 (il va y avoir
huit ans à quelques jours près) - l'intuition de l'effet dévastateur, dans ma \unsure{vi|vie}
comme dans celle de ma mère, d'une ``certaine force'' en moi. C'était la première
fois de ma vie que je consacrais une réflexion, si sommaire soit-elle, à ce
qu'avait été ma vie et surtout, mon enfance. C'était aussi le surlendemain du
jour où j'avais découvert le pouvoir de\quad méditation (**), et c'était la pre-\linebreak
mière fois depuis ce moment que je faisais usage de ce pouvoir, si longtemps
ignoré. C'est sans propos délibéré, par l'effet d'une impulsion profonde, comme
mû par un instinct très sûr, que la réflexion ce jour-là a fini par se diriger
vers mon enfance. Avec le recul seulement, je mesure à quel point c'était bien
à la source de ma vraie force, comme aussi du conflit et de la division en moi,
que m'avait porté alors un besoin profond de connaître. Pendant près de trois \unsure{an|ans}
je n'allais plus y revenir, distrait que j'étais pendant ces années par les \unsure{seul|seules}
questions ``d'ordre du jour'', sans me rendre compte que je restais à la \unsure{périphéri|périphérie}
du conflit dans ma vie, tout en me tenant obstinément éloigné du coeur même: de
cette enfance noyée de brumes, qui paraissait si infiniment lointaine...

% position: lower paragraph
Je viens de parcourir à nouveau, ``en diagonale'', les dix huit feuilles,
d'une densité exceptionnelle, de cette méditation cruciale dans ma vie. C'est
dans la nuit qui a suivi cette méditation, ou plutôt au petit matin après cette
nuit de méditation, que j'ai eu un rêve d'une force bouleversante - le premier
rêve aussi dans ma vie dont j'aie sondé le message, passionnément. Je ne me
rendais pas plus compte alors où j'allais et de ce qui était en train de se pas-\linebreak
ser, que l'avant-veille quand j'étais en train de ``découvrir la méditation''. Du-

% position: footnote rule
\figure{\begin{tikzpicture}\draw (0,0)--(3,0);\end{tikzpicture}}
% position: footnotes
(*) Voir la note ``Yang enterre yin - ou le Super père'', n° 108.

(**) Voir la section ``Désir et méditation'', n° 39.
% position: foot pencil
\margin{150}
